\chapter{Paper Chapter Template}
\label{chap:paper-template}

This chapter is a reusable template for a thesis chapter built around one paper.
Copy it to a paper-specific filename when you are ready to integrate an actual publication.

\section{How to Use This Template}

\begin{itemize}
  \item Copy this file to a descriptive name such as \texttt{paper-01-topic.tex}.
  \item Replace the chapter title with the paper topic rather than the paper title if the thesis needs a broader framing.
  \item Rewrite the opening so the chapter reads as part of the thesis, not as a standalone paper dump.
  \item Move highly paper-specific notation into this chapter only if it does not belong in the shared background chapter.
\end{itemize}

\section{Paper Metadata}

\TODO{Record the original paper title, venue, year, and author list.}

\section{Thesis Framing}

\TODO{Explain why this paper belongs in the thesis and how it connects to the chapters before and after it.}

\section{Adapted Abstract or Chapter Overview}

\TODO{Write a chapter-specific overview that matches the tone and structure of the thesis.}

\section{Main Results}

\TODO{Integrate the methods, experiments, theory, and discussion in thesis style.}

\placeholderfigure{Replace this placeholder with a figure stored in the figures/ directory.}

\begin{table}[htbp]
  \centering
  \begin{tabular}{lll}
    \toprule
    Experiment & Metric & Notes \\
    \midrule
    Placeholder A & 0.00 & Replace with real result \\
    Placeholder B & 0.00 & Replace with real result \\
    \bottomrule
  \end{tabular}
  \caption{Replace this table with a result summary, ablation, or dataset overview.}
\end{table}

\section{Adaptation Notes}

\begin{itemize}
  \item \TODO{State what material is reproduced, adapted, or newly written for the thesis.}
  \item \TODO{Record any notation changes made to align this chapter with the rest of the thesis.}
  \item \TODO{List which figures or tables should be redrawn for consistency.}
\end{itemize}

